% Bagian: first-order-logic
% Bab: introduction
% Subbagian: models-theories

\documentclass[../../../include/open-logic-section]{subfiles}

\begin{document}

\olfileid[id]{fol}{int}{mod}

\olsection{Model dan Teori}

Setelah mendefinisikan sintaksis dan semantik logika orde pertama, kita dapat
mulai menyelidiki sifat-sifat !!{structure}s dan gagasan-gagasan semantis.
Kita juga dapat mendefinisikan sistem !!{derivation} dan menyelidikinya. Bagi
suatu himpunan !!{sentence}s, kita dapat bertanya: !!{structure}s mana yang
membuat semua !!{sentence}s dalam himpunan tersebut benar? Jika diberikan
suatu himpunan !!{sentence}s~$\Gamma$, !!a{structure}~$\Struct{M}$ yang
memenuhi semuanya disebut \emph{model dari~$\Gamma$}. Kita mungkin mulai
dari~$\Gamma$ dan mencoba menemukan model-modelnya---seperti apa model-model
itu? Seberapa besar atau kecil model-model tersebut? Namun, kita juga dapat
memulai dari satu !!{structure} atau sekumpulan !!{structure}s dan bertanya:
!!{sentence}s apa yang benar di dalamnya? Adakah !!{sentence}s yang
\emph{mencirikan} !!{structure}s tersebut dalam arti bahwa kalimat-kalimat
itu benar tepat di dalam struktur-struktur tersebut? Pertanyaan semacam ini
merupakan ranah \emph{teori model}. Pertanyaan-pertanyaan tersebut juga
mendasari \emph{metode aksiomatik}: mendeskripsikan sekumpulan !!{structure}s
dengan suatu himpunan !!{sentence}s, yakni aksioma-aksioma suatu teori. Hal
ini dimungkinkan oleh pengamatan bahwa tepat !!{sentence}s yang diakibatkan
dalam logika orde pertama oleh aksioma-aksioma itulah yang benar dalam semua
model aksioma tersebut.

Sebagai contoh yang sangat sederhana, perhatikan praorder. Suatu praorder
ialah relasi~$R$ pada suatu himpunan~$A$ yang bersifat refleksif sekaligus
transitif. Suatu himpunan~$A$ dengan relasi dua-tempat $R \subseteq A \times
A$ di atasnya tepat merupakan apa yang kita perlukan untuk memberikan
!!a{structure} bagi bahasa orde pertama dengan satu simbol relasi dua-tempat~$\Obj P$:
kita menetapkan $\Domain{M} = A$ dan $\Assign{\Obj P}{M} = R$.
Karena $R$ merupakan praorder, relasi tersebut bersifat refleksif dan
transitif, dan kita dapat menemukan suatu himpunan~$\Gamma$ dari
!!{sentence}s logika orde pertama yang menyatakan hal ini:
\begin{align*}
  & \lforall[\Obj v_0][\Atom{\Obj P}{\Obj v_0,\Obj v_0}]\\
  & \lforall[\Obj v_0][\lforall[\Obj v_1][\lforall[\Obj v_2][((\Atom{\Obj P}{\Obj v_0,\Obj v_1} \land \Atom{\Obj P}{\Obj v_1,\Obj v_2}) \lif \Atom{\Obj P}{\Obj v_0,\Obj v_2})]]]
\end{align*}
Kedua !!{sentence}s ini hanyalah penyimbolan dari ``untuk setiap~$x$, $Rxx$''
($R$ bersifat refleksif) dan ``setiap kali $Rxy$ dan $Ryz$, berlaku pula
$Rxz$'' ($R$ bersifat transitif). Kita melihat bahwa
!!a{structure}~$\Struct{M}$ merupakan model dari kedua !!{sentence}s dalam~$\Gamma$
ini jika dan hanya jika $R$ (yakni, $\Assign{\Obj P}{M}$)
merupakan praorder pada~$A$ (yakni, $\Domain{M}$). Dengan kata lain,
model-model $\Gamma$ tepat merupakan semua praorder. Setiap sifat semua
praorder yang dapat diungkapkan dalam bahasa orde pertama dengan hanya~$\Obj P$
sebagai !!{predicate} (seperti refleksivitas dan transitivitas di
atas) diakibatkan oleh kedua !!{sentence}s dalam~$\Gamma$, dan sebaliknya.
Jadi, apa pun yang dapat kita buktikan mengenai model-model~$\Gamma$ telah
kita buktikan mengenai semua praorder.

Bagi setiap teori dan kelas model tertentu (seperti $\Gamma$ dan semua
praorder), akan ada pertanyaan menarik mengenai apa yang dapat dan tidak dapat
diungkapkan dalam bahasa orde pertama yang bersesuaian. Ada beberapa sifat
!!{structure}s yang menarik bagi semua bahasa dan kelas model, yakni sifat
yang menyangkut ukuran !!{domain}. Misalnya, kita selalu dapat menyatakan
bahwa !!{domain} memuat tepat $n$~!!{element}s, bagi setiap~$n \in \PosInt$.
Dengan menggunakan suatu himpunan yang memuat tak hingga banyaknya
!!{sentence}s, kita juga dapat menyatakan bahwa !!{domain} bersifat tak
hingga. Namun, kita tidak dapat menyatakan bahwa domain bersifat hingga, atau
bahwa domain tersebut !!{nonenumerable}. Hasil-hasil mengenai keterbatasan
bahasa orde pertama ini merupakan konsekuensi teorema kekompakan dan teorema
L\"owenheim--Skolem.

\end{document}
